\begin{description}
  \item[Q1.] The paper reports Table~4 qubit counts rounded to the nearest 100.
  Our recomputation shows the raw $f(n)$ values are consistently
  \emph{lower} than the paper's rounded numbers by $\sim20$--$40$ qubits at
  small $n$ (e.g.\ $f(163)=1180.4$, paper says~1200; $f(512)=3630.0$,
  paper says~3600). At $n=512$ the raw $f$ is actually \emph{above} the
  rounded value.  Is the rounding just presentation, or does the paper
  add slack for something $\varepsilon=10$ doesn't capture (measurement,
  syndrome, spare ancillas)?  A careful re-derivation of $\varepsilon$
  from Section~5.4.2 would resolve this.
  \item[Q2.] Our order-8 subgroup demo has $q=2^n$ so the QFT peaks are
  Kronecker-delta; every measurement single-handedly recovers $s$.  What
  happens at prime $q$ that does \emph{not} divide $2^n$ (the realistic
  case)?  A minimal Qiskit follow-on that picks an EC subgroup of prime
  order (say $q=11$ or $q=13$) and applies a $2^n$-QFT for $n$ big enough
  to give $\Pr[\text{success}]\!\gtrsim\!1/(2\pi^2)$ per shot would let
  us empirically measure success probability vs.\ $n{-}\log_2 q$ and
  compare with the paper's Appendix~A.2 bound.
  \item[Q3.] The paper's ``quantum halting problem'' handling (Section~5.3.4)
  produces $O(\log n)$ ``halting-time'' garbage qubits per Euclid call, and
  the analysis assumes those garbage bits remain unentangled with the
  computational registers except through the halting time itself.  Has
  that unentanglement been verified numerically for small $n$
  (say $n\!\in\!\{4,5,6\}$)?  A statevector simulation of the actual
  reversible Euclid circuit at those sizes would show whether the
  ``little bit of garbage'' really is little.
  \item[Q4.] The paper compares against Beauregard's $2n$-qubit factoring
  algorithm and concludes ``$1000$ vs $2000$ qubits'' at security-equivalent
  key sizes.  Since 2003 the RSA side has been re-analysed by
  H\"aner--Roetteler--Svore (2017), Gidney (2019), Gidney--Ekera (2021),
  etc., yielding much tighter counts.  If the same modern-vintage
  optimisations (windowed arithmetic, coset representations, deferred
  modular reduction) were applied to Proos--Zalka's ECC circuit, does
  the $1000/2000$ ratio hold, shrink, or invert?  Section~5.4.3 lists
  ``possible improvements'' but does not quantify them.
  \item[Q5.] Our reproduction treats the group-shift oracle at the level of
  a permutation on group indices ($|k\rangle\mapsto|k{+}i\bmod q\rangle$).
  The paper's Section~4 spells out the arithmetic implementation
  ($x,y\mapsto x+\lambda(y-x),\dots$) using a specific projective
  coordinate.  Are there alternative curve representations
  (\emph{e.g.}\ Edwards, Montgomery, or complete addition formulas of
  Renes--Costello--Batina, 2016) where the reversible arithmetic circuit
  is meaningfully cheaper in qubits or Toffoli count?  The paper's
  formulas assume one specific affine formula and do not survey
  alternatives.
\end{description}
